The composition of the datasets for the single- and multi-section manipulators is shown in Table~S3, Table~S4, and Table~S12. The sampled workspace of the multi-section robot is shown in Fig.~S1a, and the main reference systems in Fig.~S1b. The single-section dataset (10,000 samples) was recorded in one day by motor babbling with tendon displacements sampled uniformly in $[-45, 10]$~mm. The multi-section dataset (10,000 samples) was recorded on two days with the strategy of Algorithm~S1: for each commanded configuration, one of the three sampling modes is selected with equal probability, so that each mode accounts for about one third of the configurations; the mode was not stored for each sample. Each commanded configuration was held for about 3~s while samples were logged at 1~Hz. The forward-kinematics dataset of $N_p^3$ (10,000 samples) was recorded in three sessions over two further days. All samples of each dataset were pooled and randomly split (85/10/5\%, fixed seed). \textbf{[FILL: Data were saved in batches that were inspected during acquisition; X batches (Y samples) were discarded because ...]}

Because of this acquisition scheme, test samples can be close to training samples. For the single-section arm, 52.6\% of the test samples have a training sample within 1~mm in tendon space, but their error is not lower than that of the others (4.08 vs.\ 3.86~mm). For the multi-section arm, 26.8\% of the test samples share their configuration with a training sample; excluding them increases the tip position errors to 4.60, 14.72, and 27.99~mm, compared with 4.44, 14.01, and 26.45~mm on the full test set (Table~S14). The control targets were selected among the sampled workspace points (Section~6.3) and were therefore not excluded from training; the closed-loop experiments use new camera images, and the trajectories are not part of the datasets.
